(i) $\omega^2 = (-1)^{10}(-1)^4 = +1$ by Fact \claimref{rule2.imported-classification-centre-type-c}, and $p - q = -3 \equiv 1 \pmod 4$; $\gamma e_4 = e_0e_1e_2e_3e_4 = \omega$ by definition. (ii) The map $a \mapsto \varepsilon_+ a$ is a homomorphism because $\varepsilon_+$ is a central idempotent. Injectivity: $\varepsilon_+ a = 0$ means $a = \omega a$; but for $a \in \Cl(1,3)$ every blade of $a$ omits $e_4$ while every blade of $\omega a = \gamma e_4 a$ contains it, so $a = \omega a$ forces $a = 0$. Both $\Cl(1,3)$ and the ideal have dimension $16$ (the ideal is $\Cl^0(1,4)$ by Fact \claimref{rule2.imported-rule-2-matrix-units-forgotten}), so the map is onto. For the fifth generator, $e_4 = \gamma^{-1}\omega = -\gamma\omega$, so $\varepsilon_+ e_4 = -\gamma\varepsilon_+\omega = -\gamma\varepsilon_+$ since $\omega\varepsilon_+ = \varepsilon_+$, and $\varepsilon_- e_4 = -\gamma\varepsilon_-\omega = +\gamma\varepsilon_-$ since $\omega\varepsilon_- = -\varepsilon_-$. Fact \claimref{rule2.imported-gamma-matrices-odd-dimension} identifies the two assignments $\Gamma^4 = \mp\Gamma^0\Gamma^1\Gamma^2\Gamma^3$ as the two inequivalent representations, and $\Cl(1,3) \cong \Mat_2(\HH)$ (\claimref{std-algebra.sta-matrix}) has the one irreducible module $\HH^2$, the Dirac spinor. (iii) $\Ad_{e_4}$ negates $e_0, \ldots, e_3$ (each anticommutes with $e_4$, \claimref{std-algebra.orthogonal-anticommute}) and fixes $e_4$; the grade involution negates all five; the composite fixes $e_0, \ldots, e_3$ and negates $e_4$, which is $R$ on generators, hence on the algebra. It negates $\omega$ (one generator negated) and so exchanges $\varepsilon_\pm$; an inner automorphism fixes the centre pointwise, so $R$ is outer. (iv) is Fact \claimref{rule2.imported-rule-2-matrix-units-forgotten} with Lemma \claimref{rule2.even-subalgebra-one-dimensional-extension} for $\Cl^0(1,4) \cong \Cl(1,3)$, and $\Mat_2(\Mat_2(\HH)) = \Mat_4(\HH)$. (v) $\Ad_\nu = \hat{\cdot}$ on $A_5$ is Fact \claimref{rule2.imported-rule-2-matrix-units-forgotten}, and $\Ad_{\nu e_4} = \Ad_\nu \circ \Ad_{e_4} = \hat{\cdot} \circ \Ad_{e_4} = R$ on $A_5$ by (iii). $(\nu e_4)^2 = -\nu^2 e_4^2 = \nu^2$, which is $+1$ for $\Cl(2,4)$ and $-1$ for $\Cl(1,5)$; a bivector of square $+1$ generates a boost, of square $-1$ a rotation. (vi) is Theorem \claimref{rule2.where-obstruction-sits-commuting-square}(c), (d) at $(p,q) = (1,3)$, with Fact \claimref{rule2.imported-completion-spacetime-algebra-spinors} naming the algebra. The script checks every item: centrality and square of $\omega$, the idempotent identities, the rank $16$ of $a \mapsto \varepsilon_\pm a$ on the sixteen blades of $\Cl(1,3)$ and of each ideal, the two identities for $\varepsilon_\pm e_4$, $R = \hat{\cdot} \circ \Ad_{e_4}$ on generators, the matrix-unit relations and the trivial centre in both repairs, $\Ad_\nu$ and $\Ad_{\nu e_4}$ on generators in both, the two squares of $\nu e_4$, the ladder relations in $\Cl(2,4)$, the definiteness in $\Cl(1,5)$, the residues, and the Rule 1 corner $P\,\Cl(2,4)\,P$ of dimension $16$ commuting with $\Cl(1,3)$.
